% ============================================================
%  Глава 6 — Одометрия: кинематика дифференциального привода
%  Файлы: ros_ws/.../motor_node.cpp  и  pi_nodes/nodes/motor_node.py
% ============================================================
\chapter{Одометрия: кинематика дифференциального привода}
\label{ch:odometry}

\begin{filebox}[title={Файлы исходного кода}]
\filepath{ros\_ws/src/robot\_pkg\_cpp/src/motor\_node.cpp} \quad (C++, строки 244--302)\\
\filepath{pi\_nodes/nodes/motor\_node.py} \quad (Python)
\end{filebox}

\section{Модель дифференциального привода}

Гусеничный робот управляется двумя независимыми бортами (левый/правый).
Линейная и угловая скорости центра масс:

\begin{equation}
  v_{\text{linear}} \in [-v_{\max}, v_{\max}], \quad
  \omega_{\text{angular}} \in [-\omega_{\max}, \omega_{\max}]
  \label{eq:velocity_limits}
\end{equation}

\section{Микшер дифференциального привода}

\begin{filebox}[title={\filepath{motor\_node.cpp}, строки 251--266 --- controlLoop()}]
\begin{lstlisting}[style=cppstyle, numbers=left, firstnumber=251]
// Convert m/s -> throttle in [-1, 1]
double lin_t = (max_lin_ > 0.0) ? (linear_  / max_lin_) : 0.0;
double ang_t = (max_ang_ > 0.0) ? (angular_ / max_ang_) : 0.0;

// Differential-drive mixer
double left  = lin_t + ang_t;
double right = lin_t - ang_t;

// Normalise: keep ratio, clamp to [-1, 1]
double max_val = std::max({std::abs(left), std::abs(right), 1.0});
if (max_val > 1.0) { left /= max_val; right /= max_val; }
\end{lstlisting}
\end{filebox}

\textbf{Формулы микшера} (нормированные команды $\in [-1, 1]$):

\begin{equation}
  \text{left} = \ell_{\text{lin}} + \ell_{\text{ang}}, \quad
  \text{right} = \ell_{\text{lin}} - \ell_{\text{ang}}
  \label{eq:diff_mixer}
\end{equation}

Нормировка для сохранения соотношения:

\begin{equation}
  m = \max\!\bigl(1,\, |\text{left}|,\, |\text{right}|\bigr), \quad
  \text{left} \leftarrow \frac{\text{left}}{m}, \quad
  \text{right} \leftarrow \frac{\text{right}}{m}
  \label{eq:normalize}
\end{equation}

\section{Интегрирование одометрии (кинематические уравнения)}

\begin{filebox}[title={\filepath{motor\_node.cpp}, строки 268--273}]
\begin{lstlisting}[style=cppstyle, numbers=left, firstnumber=268]
// Open-loop odometry integration
x_     += linear_  * std::cos(theta_) * dt;
y_     += linear_  * std::sin(theta_) * dt;
theta_ += angular_ * dt;
\end{lstlisting}
\end{filebox}

\textbf{Уравнения кинематики одноколейного робота}
(дискретная интеграция методом Эйлера вперёд):

\begin{align}
  x_{k+1} &= x_k + v \cos(\theta_k) \cdot \Delta t \label{eq:odom_x}\\
  y_{k+1} &= y_k + v \sin(\theta_k) \cdot \Delta t \label{eq:odom_y}\\
  \theta_{k+1} &= \theta_k + \omega \cdot \Delta t \label{eq:odom_theta}
\end{align}

Это модель \textbf{unicycle} (одноколейная модель): скорость $v$ направлена
вдоль текущего курса $\theta$.

\begin{notebox}
\textbf{Открытоконтурная одометрия}: модель не учитывает проскальзывание
гусениц. Именно для компенсации этой ошибки и используется Velocity EKF
(глава~\ref{ch:velocity_ekf}), который добавляет коррекцию от акселерометра.
\end{notebox}

\section{Кватернионное кодирование угла курса}

\begin{filebox}[title={\filepath{motor\_node.cpp}, строки 276--291 --- publishOdom()}]
\begin{lstlisting}[style=cppstyle, numbers=left, firstnumber=276]
void publishOdom(const rclcpp::Time & stamp)
{
  double sin_h = std::sin(theta_ / 2.0);
  double cos_h = std::cos(theta_ / 2.0);

  nav_msgs::msg::Odometry odom;
  odom.pose.pose.orientation.z = sin_h;
  odom.pose.pose.orientation.w = cos_h;
  odom.twist.twist.linear.x    = linear_;
  odom.twist.twist.angular.z   = angular_;
  odom_pub_->publish(odom);
}
\end{lstlisting}
\end{filebox}

Поворот в 3D вокруг вертикальной оси $Z$ кодируется кватернионом:

\begin{equation}
  \vect{q} = \bigl(\underbrace{0}_{q_w\text{'s missing part}},\,
  \underbrace{0}_{q_x},\,
  \underbrace{0}_{q_y},\,
  \underbrace{\sin(\theta/2)}_{q_z},\,
  \underbrace{\cos(\theta/2)}_{q_w}\bigr)
  \label{eq:quaternion_yaw}
\end{equation}

Для чистого вращения вокруг $Z$: $q_x = q_y = 0$, поэтому достаточно
хранить только $q_z$ и $q_w$ (что делает ROS2 сообщение \texttt{Odometry}).

\section{Сводная диаграмма}

\begin{center}
\begin{tikzpicture}[
  block/.style={rectangle, draw=blue!40, fill=blue!7, rounded corners,
                minimum width=3.5cm, minimum height=0.9cm,
                font=\small, align=center},
  arrow/.style={-{Stealth}, blue!60, thick}
]
  \node[block] (cmd) {Команда\\$(v_{\text{lin}},\, \omega_{\text{ang}})$};
  \node[block, right=2.5cm of cmd] (mix) {Микшер\\$L = \ell_{lin}+\ell_{ang}$\\$R = \ell_{lin}-\ell_{ang}$};
  \node[block, right=2.5cm of mix] (mot) {Моторы\\ШИМ $\in [-1, 1]$};
  \node[block, below=1.5cm of cmd] (odom) {Одометрия\\$x,y,\theta$};

  \draw[arrow] (cmd) -- (mix);
  \draw[arrow] (mix) -- (mot);
  \draw[arrow] (cmd) -- node[left, font=\tiny]{$v, \omega, \Delta t$} (odom);
\end{tikzpicture}
\end{center}

\section{Теоретическое обоснование: кинематика и динамика}

\subsection{Неголономная модель робота}

Гусеничный робот подчиняется \textbf{неголономному ограничению} ---
не может двигаться боком. В терминах ТАУ это выражается как:

\begin{equation}
  \dot{x}\sin\theta - \dot{y}\cos\theta = 0
  \label{eq:nonholonomic}
\end{equation}

Это ограничение не интегрируется в голономную связь вида $f(x, y, \theta) = 0$,
что делает систему \textbf{недоуправляемой} --- размерность конфигурационного
пространства $(x, y, \theta) = 3$ больше числа управлений $(v, \omega) = 2$.

\subsection{Модель unicycle в пространстве состояний}

Кинематическая модель~\eqref{eq:odom_x}--\eqref{eq:odom_theta} ---
это \textbf{нелинейная} система $\dot{\vect{x}} = \vect{f}(\vect{x}, \vect{u})$:

\begin{equation}
  \begin{pmatrix} \dot{x} \\ \dot{y} \\ \dot{\theta} \end{pmatrix}
  = \begin{pmatrix} \cos\theta & 0 \\ \sin\theta & 0 \\ 0 & 1 \end{pmatrix}
  \begin{pmatrix} v \\ \omega \end{pmatrix}
  \label{eq:unicycle_ss}
\end{equation}

Якобиан по состоянию (для EKF):

\begin{equation}
  \mat{F} = \frac{\partial \vect{f}}{\partial \vect{x}} =
  \begin{pmatrix}
    1 & 0 & -v\sin\theta\,\Delta t \\
    0 & 1 & v\cos\theta\,\Delta t \\
    0 & 0 & 1
  \end{pmatrix}
  \label{eq:unicycle_jacobian}
\end{equation}

\subsection{Связь с моделью двигателя}

Каждый борт робота приводится двигателем постоянного тока
(гл.~\ref{ch:theory}, раздел~\ref{sec:dc_motor}). Полная
динамическая модель включает электромеханику:

\begin{equation}
  T_a^2\ddot{\omega}_{\text{motor}} + T_y\dot{\omega}_{\text{motor}} + \omega_{\text{motor}} = k\,u_{\text{PWM}}
  \label{eq:motor_dynamics}
\end{equation}

Однако в проекте используется \textbf{кинематическая модель}
(без учёта инерции), что оправдано малым отношением
$T_a/T_{\text{control}} \ll 1$ (электрическая постоянная времени
много меньше периода управления $50\,\text{мс}$).

\subsection{Погрешность метода Эйлера}

Интегрирование методом Эйлера вперёд имеет погрешность $O(\Delta t^2)$
на шаг и $O(\Delta t)$ глобально. Для периода $\Delta t = 0.05\,\text{с}$
и $v_{\max} = 0.5\,\text{м/с}$ максимальная ошибка позиции за шаг:

\begin{equation}
  \epsilon \leq \frac{1}{2}|\ddot{x}|\Delta t^2 \approx \frac{v\omega\Delta t^2}{2}
  \leq \frac{0.5 \cdot 2.0 \cdot 0.0025}{2} = 1.25\,\text{мм}
  \label{eq:euler_error}
\end{equation}

Эта ошибка компенсируется Velocity EKF (гл.~\ref{ch:velocity_ekf}).

\section{Параметры}

\begin{center}
\begin{tabular}{lll}
\toprule
\textbf{Параметр} & \textbf{Значение} & \textbf{Описание}\\
\midrule
\texttt{max\_lin\_} & $0.5\,\text{м/с}$ & Максимальная линейная скорость\\
\texttt{max\_ang\_} & $2.0\,\text{рад/с}$ & Максимальная угловая скорость\\
Частота & $20\,\text{Гц}$ & Период управляющего цикла ($\Delta t = 0.05\,\text{с}$)\\
\bottomrule
\end{tabular}
\end{center}
